\documentclass[11pt]{article}
\usepackage[margin=1in]{geometry}
\usepackage{amsmath,amssymb}
\usepackage{hyperref}

\title{Literature-implied answer: high-dimensional Krivine schemes are asymptotically optimal}
\author{Run fa\_banach\_001, agent lane 17}
\date{2026-07-03}

\begin{document}
\maketitle
\sloppy

\section*{Source Question}

The source paper is Braverman--Makarychev--Makarychev--Naor,
\emph{The Grothendieck constant is strictly smaller than Krivine's bound},
arXiv:1103.6161 \cite{BMMN2011}. After disproving Krivine's conjecture and
Konig's intermediate conjecture, the authors ask whether one can analytically
describe the high-dimensional maximizers
\[
 f_{\max}^{(n)},g_{\max}^{(n)}:\mathbb R^n\to\{-1,1\}
\]
of Konig's bilinear form, and whether these maximizers form an alternating
Krivine rounding scheme. They note that such a positive answer would yield a
formula recovering \(K_G\) from high-dimensional operator norms.

\section*{Supporting Theorem}

Naor--Regev, \emph{Krivine schemes are optimal}, arXiv:1205.6415
\cite{NR2012}, proves the following theorem: for every \(k\in\mathbb N\)
there is a \(k\)-dimensional oblivious Krivine scheme of quality
\[
 (1+O(1/k))K_G .
\]
Their introduction explicitly frames this as showing that oblivious Krivine
rounding schemes capture \(K_G\) exactly in the asymptotic sense: to understand
the Grothendieck constant it suffices, up to a factor tending to one, to focus
on partitions of Euclidean space used in Krivine schemes.

\section*{Conclusion}

This gives a literature-implied answer to the broad high-dimensional direction
opened in arXiv:1103.6161: high-dimensional Krivine schemes are asymptotically
optimal and lose only a \(1+O(1/k)\) factor.

The scope is partial. The later theorem does not identify the actual
maximizers of Konig's bilinear form, does not give an analytic formula for the
``tiger partition'' \(f_\infty\), and does not prove convergence of the
iteration described in the source paper. Those analytic/maximizer questions
should remain open in the run memory.

\begin{thebibliography}{9}
\bibitem{BMMN2011}
M. Braverman, K. Makarychev, Y. Makarychev, and A. Naor,
\emph{The Grothendieck constant is strictly smaller than Krivine's bound},
arXiv:1103.6161, 2011.

\bibitem{NR2012}
A. Naor and O. Regev,
\emph{Krivine schemes are optimal}, arXiv:1205.6415, 2012.
\end{thebibliography}

\end{document}
